% !TeX root = ../main.tex

\chapter{总结与展望}

\section{工作总结}

本文基于安卓和龙芯GPU，设计并实现了一个安卓图形栈，填补了LoongArch架构下安卓图形系统的空缺，可以运行在龙芯常见的硬件环境下。由于GPU的软硬件环境较为复杂，本文仅涉及
部分关键模块及库的构建。主要工作包括：

1.在龙芯安卓系统上实现了系统依赖的两个关键功能硬件混合渲染和图形缓存分配。
硬件混合渲染模块作为SurfaceFlinger的服务端，分为DRM管理模块、HWC2接口模块、合成器模块、缓冲区信息适配模块以及后端管理模块五个子模块，
负责显示资源管理、合成计划创建、验证与显示、缓冲区信息适配等核心功能。
而图形缓存分配模块承担Gralloc功能，基于GBM接口标准实现了基于龙芯LG显卡的缓冲区创建、导入、映射、同步等核心功能，
并基于Gralloc4标准实现了安卓Gralloc HAL中Allocate和Mapper的接口功能。

2.为兼容硬件底层功能接口，以及兼容多个内核版本的DRM接口。针对这些问题，本文基于Linux DRM框架，设计了17类IOCTL命令码，
覆盖显存管理（GEM）、命令流控制（CS）、虚拟地址映射（VA）等核心功能，为从用户空间高效提交命令到GPU处理器提供支持。
针对Linux内核版本差异（4.19至5.10），提出动态编译检测方案，通过条件编译实现接口兼容性，解决了闭源驱动的多版本适配问题。

% 1.对安卓图形系统进行了需求分析。本文从关键服务分析、架构差异分析以及安卓兼容性三个方面导入需求，并对两个主要模块，图形缓存分配模块和硬件混合渲染模块进行需求分析。
% 随后介绍了图形系统的整体架构，以及功能模块的设计。

% 2.详细设计并实现了几个关键模块。内核模块实现基于Linux DRM框架，设计了17类IOCTL命令码，覆盖显存管理（GEM）、命令流控制（CS）、虚拟地址映射（VA）等核心功能。
% 针对Linux内核版本差异（4.19至5.10），提出动态编译检测方案，通过条件编译实现接口兼容性，解决了闭源驱动的多版本适配问题。
% 图形缓存分配器基于DMA-BUF框架，采用缓冲区复用与零拷贝技术，支持显存与系统内存的动态分配策略，优化了图形缓冲区的分配效率。
% 通过HIDL接口实现绑定式服务，确保跨进程共享的安全性。
% 硬件混合渲染器结合DRM/KMS子系统，实现了显示资源管理、合成计划创建、验证与显示、缓冲区信息适配等核心功能。

最后在系统添加基于Mesa框架的龙芯OpenGL ES2.0支持，定制系统镜像。
通过功能测试以及非功能测试并对各模块进行充分验证。
实验结果表明，系统在国产LoongArch平台上实现Android图形栈的完整运行并满足设计目标。

% 本文在龙芯平台上实现Android图形系统的全栈适配，基于龙芯显卡实现了gralloc与HWC模块，并构建OpenGL ES和Tile访存优化，实现了龙芯LG显卡与安卓框架的有效协同，
% 为后续基于该平台的开发和研究奠定了技术基础。

\section{未来展望}

GPU所涉及的软件支持很复杂，包括编译器后端生成、图形编程接口的支持等，本文只涉及部分模块和内容。同样安卓系统也作为一个即其庞大而复杂的系统，其图形系统更是作为最复杂的系统之一，
从工程实现和调试的角度来看都有不小的难度。本文所做的工作仅涉及部分模块的实现以及构建，接下来工作可以分成几个方向：

1.系统及硬件升级：当前方案仅适配Android S版本与龙芯3A5000平台，并且为了强行支持龙芯硬件特性对安卓系统镜像进行了许多定制，导致ART运行以及zygote等安卓APP无法运行，
在测试上也带来了一定的困难，未来需扩展至更高版本安卓（如Android 15）支持更多的硬件特性及新一代龙芯硬件（如LG200），并针对Vulkan API支持等场景进行性能优化。
从而实现在实际场景下的验证，如3D游戏、视频渲下的性能评估，运行一些常见的JAVA性能测试工具，测试帧率、功耗等关键指标。
而LG110作为一个过渡性、研究性的GPU，是龙芯制造国产GPU道路上一个尝试性的产品，其本身具有许多瑕疵，不适合继续深入做研究，应尽快过渡至LG200。

2.该方案目前的重点是搭建一个可运行的图形环境，并探索相应的软件依赖方案，其实现上仍有许多不足。在完成对功能的完善后可以考虑更多定制化的方案进行优化，以实现更高效的图形环境。

